\chapter{Layer 2: Canonical Symbolic IR}
\label{ch:ir}

\section{Purpose and Contract}
\label{sec:ir-purpose}

Layer~2 receives the parsed AST and emits a \emph{canonical} S-expression. The
inbound contract is ``a typed \texttt{Expr}''; the outbound contract is ``a
normalized representation such that semantically equal expressions have
syntactically identical IR.'' Canonicalization is what makes later
equality, proof, and caching meaningful.

\section{The IR Datatype}
\label{sec:ir-ast}

The reference IR (\texttt{layers/sexpr/src/AST.hs}) is an S-expression
datatype:
\begin{lstlisting}[language=Haskell,caption={S-expression AST (AST.hs).},label=lst:ir-ast]
data SExpr
  = SInt Integer
  | SVar String
  | SApp String [SExpr]
  deriving (Show, Eq)
\end{lstlisting}
This mirrors the Layer~1 \texttt{Expr} but collapses the distinction between
lambdas and primitives into a uniform application node \texttt{SApp name
args}. Uniformity is deliberate: it makes the IR a single recursive type that
every downstream tool (Lean bridge, normalizer, FOL checker) can consume
without case-splits on language-specific constructors.

\section{Normalization}
\label{sec:ir-norm}

Canonicalization applies, at minimum:
\begin{enumerate}
  \item \textbf{Constant folding} of integer literals.
  \item \textbf{Commutativity-aware ordering} of associative commutative
        operators so that $a+b$ and $b+a$ share a form.
  \item \textbf{Flattening} of nested same-tag nodes (e.g., $(a+b)+c
        \mapsto a+b+c$).
\end{enumerate}
The C{--} kernel mirrors this at Layer~5 with \texttt{normalize\_with\_cert},
which ``applies associativity/commutativity, expands polynomials, and
generates a receipt seal'' (\texttt{verified\_kernel.cm}). The two
normalizers operate at different layers but share the same contract: equal
expressions in, equal expressions out.

\section{From AST to IR}
\label{sec:ir-conv}

The conversion is a straightforward fold over \texttt{Expr}:
\begin{itemize}
  \item \texttt{IntVal i} $\mapsto$ \texttt{SInt i}
  \item \texttt{Var v} $\mapsto$ \texttt{SVar v}
  \item \texttt{App f xs} $\mapsto$ \texttt{SApp (render f) (map conv xs)}
  \item \texttt{Lambda [..] e} $\mapsto$ \texttt{SApp "lambda" [conv e]}
  \item \texttt{Prim p xs} $\mapsto$ \texttt{SApp p (map conv xs)}
\end{itemize}
where \texttt{render} prints the function expression into a stable string key.

\section{Why a Canonical IR Matters}
\label{sec:ir-why}

Three downstream benefits justify the layer:
\begin{enumerate}
  \item \textbf{Proof stability.} Lean obligations generated from a canonical
        form are stable across syntactic variation, so re-running the pipeline
        on equivalent inputs yields comparable goals.
  \item \textbf{Cache hits.} Two queries that are equal after normalization
        can share a verified kernel result and a WORM receipt.
  \item \textbf{Simpler analysis.} A single recursive type reduces the
        surface area of the static-analysis and FOL layers.
\end{enumerate}
The IR is the spine of the entire pipeline: every later layer either consumes
it or produces an artifact derived from it.
